\documentclass[10pt,a4paper]{amsart}
\usepackage[margin=19mm]{geometry}
\usepackage{amsmath,amssymb,mathtools}
\usepackage[scr=rsfs,cal=euler]{mathalfa}
\usepackage{xfrac}
\usepackage[unicode,hidelinks,bookmarks=false]{hyperref}
\newcommand{\ie}{i.e.\ }
\newcommand{\cf}{cf.\ }
\newcommand{\resp}{respectively\ }
\newcommand{\Subsection}{\subsection}
\providecommand{\nobreakdash}{\nobreak\hbox{-}}
\setlength{\emergencystretch}{2em}
\multlinegap0pt
\usepackage{bm}
\usepackage{amssymb}
\newtheorem{lem}{Lemma}
\newtheorem{prop}{Proposition}
\title{The derived set of multiple star polylogarithms}
\author{Jiangtao Li}
\date{Source: arXiv:2609.03653v1 (CC BY 4.0)}
\begin{document}
\maketitle

\noindent\emph{Source and licence.} The derived set of multiple star polylogarithms. Version arXiv:2609.03653v1 (CC BY 4.0), distributed by arXiv under the Creative Commons Attribution 4.0 International licence, http://creativecommons.org/licenses/by/4.0/, as recorded in the arXiv licence metadata for this exact version and shown on the abstract page https://arxiv.org/abs/2609.03653v1. Full licence notice: \texttt{licenses/arxiv-2609.03653-CC-BY-4.0.txt}.\par\smallskip

\noindent\textit{Editorial scope.} Counted unit: the article's prop logderivative (source0.tex lines 951--957). The article's complete original proof of it is delivered with it (source0.tex lines 959--1049, one \texttt{proof} environment of 91 lines). No mathematics is written, repaired or re-derived: the statement and the proof are the article's own lines, in the article's own notation, and no retained line is substituted or re-flowed.  Retained uncounted as the source-local context the proof needs: lines 590--593; lines 595--598; lines 760--773; lines 797--803; lines 829--835; lines 937--940; lines 944--949.  Nothing was omitted from any retained span, so the package carries no omission record.  No retained line contains a citation, so the wrapper carries no bibliography and no bibliography entry.  No retained line contains a figure environment, an image-inclusion command or drawing code, so no image is delivered and the package declares no asset dependency.  Editorial formatting only, in addition: an AMS \texttt{amsart} preamble that restates the article's own declarations and macros used by the retained lines, namely the environments \texttt{lem}, \texttt{prop} on separate counters, together with the standard packages those lines need.  The retained environments print in this document as Lemma 1, Proposition 1 in order of appearance, whereas the article prints the same items with the numbering of its own full document.  The complete native source of the article as distributed in the arXiv e-print for this exact version is bundled unchanged as \texttt{sources/209311/source0.tex}.
\begin{equation}\label{ArecM}
A_{(p,{\boldsymbol\tau})}(n)
=\frac1{n+1}A_{(p-1,{\boldsymbol\tau})}(n).
\end{equation}

\begin{equation}\label{ArecC}
A_{(1,{\boldsymbol\tau})}(n)
=\frac1{n+1}\sum_{j=0}^{n}A_{\boldsymbol\tau}(j).
\end{equation}

\begin{lem}\label{quotientlemma}
Suppose that
\[
f(z)=\int_0^1\frac{\phi(t)dt}{1-tz},
\qquad
g(z)=c+\int_0^1\frac{\psi(t)dt}{1-tz},
\]
where $c\geq 0$, $f,g\in\mathfrak{T}$, $\phi,
\psi>0$ on $(0,1)$, and $\phi/\psi$ is nondecreasing.  If $c=0$, assume in addition that
\begin{equation}\label{psiintegrable}
\int_0^1\frac{\psi(t)}t\,dt<+\infty.
\end{equation}
Then $f/g\in\mathfrak{T}$.
\end{lem}

\begin{lem}\label{involution}
If $h\in\mathfrak{T}$, then
\begin{equation}\label{Jdef}
(\mathcal{J}h)(z)=\frac1{(1-z)h(z)}
\end{equation}
belongs to $\mathfrak{T}$.
\end{lem}

\begin{prop}\label{strictMLR}
If ${\bf k}\succ{\boldsymbol\ell}$, then
\begin{equation}\label{MLRratio}
t\mapsto\frac{\sigma_{\bf k}(t)}{\sigma_{\boldsymbol\ell}(t)}
\end{equation}
is strictly increasing on $(0,1)$.
\end{prop}

\begin{equation}\label{denintegrable}
\int_0^1\frac{t\sigma_{\bf k}(t)/A_{\bf k}(1)}t\,dt
=\frac1{A_{\bf k}(1)}<+\infty.
\end{equation}

\begin{equation}\label{Fminusinfty}
F_{\bf k}(-x)
=-\int_0^1\frac{xt}{1+xt}\sigma_{\bf k}(t)\,dt
\rightarrow -1
\qquad(x\rightarrow+\infty).
\end{equation}

\begin{prop}\label{logderivative}
For every nonempty finite index ${\bf k}$, the function
\begin{equation}\label{Pdef}
P_{\bf k}(z)=\frac{zF_{\bf k}'(z)}{F_{\bf k}(z)}
\end{equation}
has a removable singularity at $0$, belongs to $\mathfrak{T}$ after that singularity is filled in, and is not constant.
\end{prop}

\begin{proof}
Write ${\bf k}=(p,{\boldsymbol\tau})$. Assume first that $p\geq 2$, and put ${\bf k}^{-}=(p-1,{\boldsymbol\tau})$.  Then ${\bf k}^{-}\succ{\bf k}$.  Proposition \ref{strictMLR}, Lemma \ref{quotientlemma}, and \eqref{denintegrable} give
\begin{equation}\label{qminus}
q(z)=\frac{\widehat F_{{\bf k}^{-}}(z)}{\widehat F_{\bf k}(z)}
=\frac12\frac{F_{{\bf k}^{-}}(z)}{F_{\bf k}(z)}
\in\mathfrak{T},
\end{equation}
because \eqref{ArecM} gives
$A_{{\bf k}^{-}}(1)=2A_{\bf k}(1)$.  By \eqref{Fminusinfty}, $q(-x)\rightarrow1/2$, so the probability measure representing $q$ has an atom of mass $1/2$ at $0$.  The coefficient identity \eqref{ArecM} is equivalent to
\[
F_{{\bf k}^{-}}=F_{\bf k}+zF_{\bf k}'.
\]
Consequently,
\[
P_{\bf k}=\frac{F_{{\bf k}^{-}}}{F_{\bf k}}-1
=2q-1
=2\int_{(0,1]}\frac{d\nu(t)}{1-tz}\in\mathfrak{T},
\]
because the measure on the right has total mass $1$.

Now assume that $p=1$.  Put
\[
G_{\boldsymbol\tau}(z)=z+F_{\boldsymbol\tau}(z),
\qquad
h(z)=2\frac{F_{(1,{\boldsymbol\tau})}(z)}{G_{\boldsymbol\tau}(z)}.
\]
When ${\boldsymbol\tau}=\varnothing$, one has
$h=2F_{(1)}/z=\widehat F_{(1)}\in\mathfrak{T}$.  Suppose that ${\boldsymbol\tau}\neq\varnothing$ and set
\[
S=1+A_{\boldsymbol\tau}(1).
\]
The case $n=1$ of \eqref{ArecC} gives
\begin{equation}\label{halfcoefficient}
2A_{(1,{\boldsymbol\tau})}(1)=S.
\end{equation}
Therefore
\[
f(z)=\frac{2F_{(1,{\boldsymbol\tau})}(z)}{Sz}
=\int_0^1\frac{2t\sigma_{(1,{\boldsymbol\tau})}(t)/S}{1-tz}\,dt
\in\mathfrak{T},
\]
\[
g(z)=\frac{G_{\boldsymbol\tau}(z)}{Sz}
=\frac1S+\int_0^1\frac{t\sigma_{\boldsymbol\tau}(t)/S}{1-tz}\,dt
\in\mathfrak{T}.
\]
The representing measure of $g$ has a positive atom at $0$.  Moreover,
$(1,{\boldsymbol\tau})\succ{\boldsymbol\tau}$.  To see this, if all components of ${\boldsymbol\tau}$ are $1$, then ${\boldsymbol\tau}$ is a proper prefix of $(1,{\boldsymbol\tau})$; otherwise, at the first component of ${\boldsymbol\tau}$ greater than $1$, the corresponding component of $(1,{\boldsymbol\tau})$ is $1$.  Proposition \ref{strictMLR} and Lemma \ref{quotientlemma} now give
\[
h=\frac fg\in\mathfrak{T}.
\]

By Lemma \ref{involution}, $\mathcal{J}h\in\mathfrak{T}$.  Equation \eqref{Fminusinfty} gives
\[
h(-x)\sim\frac2x,
\qquad
(\mathcal{J}h)(-x)=\frac1{(1+x)h(-x)}\rightarrow\frac12.
\]
Thus the representing measure of $\mathcal{J}h$ has an atom of mass $1/2$ at $0$.

For ${\boldsymbol\tau}\neq\varnothing$, \eqref{ArecC} gives
\[
(n+1)A_{(1,{\boldsymbol\tau})}(n)
=1+\sum_{j=1}^{n}A_{\boldsymbol\tau}(j),
\]
and the same formula, with the sum absent, holds for
${\boldsymbol\tau}=\varnothing$.  Hence
\begin{equation}\label{firstoneidentity}
F_{(1,{\boldsymbol\tau})}+zF_{(1,{\boldsymbol\tau})}'
=\frac{z+F_{\boldsymbol\tau}}{1-z}
=\frac{G_{\boldsymbol\tau}}{1-z}.
\end{equation}
It follows that
\[
P_{(1,{\boldsymbol\tau})}+1
=\frac{G_{\boldsymbol\tau}}{(1-z)F_{(1,{\boldsymbol\tau})}}
=2\mathcal{J}h.
\]
Removing the atom of mass $1/2$ at $0$ shows that
$P_{(1,{\boldsymbol\tau})}=2\mathcal{J}h-1\in\mathfrak{T}$.

Finally, if
\[
F_{\bf k}(z)=A_{\bf k}(1)z+A_{\bf k}(2)z^2+\cdots,
\]
then
\[
P_{\bf k}(z)=1+\frac{A_{\bf k}(2)}{A_{\bf k}(1)}z+O(z^2).
\]
Since $A_{\bf k}(2)>0$, the function is not constant.
\end{proof}

\end{document}
